\documentclass[11pt,a4paper]{article}
\usepackage[utf8]{inputenc}
\usepackage[T1]{fontenc}
\usepackage{lmodern}
\usepackage{geometry}
\usepackage{xcolor}
\usepackage[colorlinks=true,linkcolor=blue,urlcolor=blue]{hyperref}
\geometry{margin=2.4cm}
\setlength{\emergencystretch}{2em}
\setlength{\parskip}{0.45em}
\setlength{\parindent}{0pt}
\usepackage{titlesec}
\titlespacing*{\section}{0pt}{1.3ex plus 1ex minus .2ex}{0.6ex}

\title{\bfseries Response to the Deputy Editor's Pre-Review Comments\\[6pt]
\normalsize Manuscript: \textit{Spatio-temporal wall-actuation modes and
cycle-mean viscous dissipation in a lid-driven cavity}}

\author{Beniaiche Ahmed (corresponding author), Larabi Abderrahim, Belkadi Mbarek\\
Laboratory of Fluid Mechanics, Ecole Militaire Polytechnique, Algiers, Algeria}

\date{}

\begin{document}
\maketitle
\vspace{-2.5em}

\noindent\textbf{To:} The Deputy Editor, \textit{Physics of Fluids}.\\
\textbf{Re:} Revision requested before external review. The revised manuscript
and this response are submitted through the PXP site by 21 November 2026.

\section*{Opening statement}

We thank the Deputy Editor for confirming that the chosen article type is
appropriate and for the pre-review comments. The manuscript has been revised to
meet them. This letter answers each comment in the order it was given. The
changes themselves are visible in the marked-up manuscript and are itemized,
comment by comment, in the point-by-point response transmitted with the
revision; they are not repeated here.

\section*{Comment 1}

\noindent\textcolor{blue}{``Physics of Fluids publishes original theoretical,
computational, and experimental work that advances the physical understanding
of fluids, meaning the mechanisms and principles that explain how and why
fluids behave as they do. As written, the manuscript is framed as a
reproducibility analysis of a physics informed neural network controller for
the lid driven cavity. Before it can be considered, the paper needs to be
reorganized so that the fluid physics is its central contribution.''}

\textbf{Response.} The manuscript has been reorganized around the fluid
physics. It now studies how the spatial and temporal structure of a wall
actuation sets the viscous dissipation of the cavity flow it drives, and the
physics-informed network appears only as the tool that generates the candidate
actuations, which an independent multiple-relaxation-time lattice-Boltzmann
solver then evaluates. The identification campaigns and the reproducibility
audit no longer frame the paper.

\section*{Comment 2}

\noindent\textcolor{blue}{``The physics to bring to the front is how the spatio
temporal structure of the wall actuation sets the cycle-mean viscous
dissipation of the cavity flow. The neural network branch selection, the
regularization dependent state recovery, and the reproducibility audit can
remain as supporting evidence, but they should not frame the work.''}

\textbf{Response.} The cycle-mean dissipation is now the quantity the paper
compares, at a common quadratic actuation norm, for the stationary modes and
the three identified branches. The distinction between the cycle-mean
dissipation and the dissipation of the time-mean field, and the partition of
the cycle-mean cost between mean field and fluctuations, are reported for each
branch. The branch selection, the regularization-dependent state recovery and
the reproducibility audit remain as supporting evidence in the supplementary
material.

\section*{Comment 3}

\noindent\textcolor{blue}{``The introduction and reference list should be
repositioned toward the fluid physics literature on active control of the lid
driven cavity, viscous dissipation and energy budgets in actuated cavity
flows, and modal structure and state selection in driven-cavity flow, so that
the contribution is set against fluid flow.''}

\textbf{Response.} The introduction and the reference list are now positioned
in that literature: the cavity benchmarks, the active-control and
cost-function work, the dissipation and energy-budget studies, and the modal
structure and state-selection results. The antecedent identification study by
one of us is cited where its results are used and is not restated.

\section*{Comment 4}

\noindent\textcolor{blue}{``Please also reduce the number of figures to maximum
of 15.''}

\textbf{Response.} The revised manuscript carries eight figures and four
tables. The figures of the identification campaigns are no longer in the main
text; those that support the comparison are in the supplementary material.

\section*{Comment 5}

\noindent\textcolor{blue}{``Please revise the title to lead with the actuation
mode and dissipation physics.''}

\textbf{Response.} The title is now \textit{Spatio-temporal wall-actuation
modes and cycle-mean viscous dissipation in a lid-driven cavity}, which leads
with the actuation structure and the dissipation cost.

\section*{Comment 6}

\noindent\textcolor{blue}{``Please revise your manuscript and submit a
detailed, point-by-point response to the comments. Each response should
include the original comment followed by a description of the corresponding
revision made (or a justification if no change was made). The revised
manuscript and response are due by 21-Nov-2026 via the PXP web site. The
Editors will then make a decision as to the next step in the review
process.''}

\textbf{Response.} A detailed point-by-point response is submitted with the
revision. It reproduces every comment of the Deputy Editor and of the
reviewers, each followed by the corresponding revision or by the justification
where no change was made. The revised manuscript and the response are
submitted through the PXP site by 21 November 2026.

\section*{Comment 7}

\noindent\textcolor{blue}{``When submitting your revision, please upload two
versions of your manuscript: a clean version labeled `Article File' and a
marked-up version labeled `Marked-Up Manuscript,' which should clearly show
all changes made from the previous submission.''}

\textbf{Response.} Both versions are uploaded: the clean manuscript as the
Article File and the marked-up manuscript as the Marked-Up Manuscript, the
latter showing every change against the previous submission. The same pair is
provided for the supplementary material.

\section*{Comment 8}

\noindent\textcolor{blue}{``Please add a section heading in the manuscript
labeled `Supplementary Material' and include a brief description of the
supplemental material. It should be placed after the Conclusion and before the
Acknowledgment. AIP Publishing will insert a link that will allow successful
linking between the published paper and the associated supplementary material.
Supplementary Material source files must have a Word version in addition to the
\LaTeX{} version.''}

\textbf{Response.} The revised manuscript now carries an unnumbered
\textit{Supplementary Material} section between the Conclusions and the
Acknowledgments, where the publisher's link is to be inserted. Its single
paragraph says which records the supplement holds and quotes its extent
(Secs.~S1--S5, Tables~S1--S19, Figs.~S1--S2); the certification script reads those
three numbers back from the compiled supplement and fails if the paragraph
disagrees with them or if the section leaves its required position. The
supplementary material is submitted as a Word source beside its \LaTeX{} source:
the Word file is generated from the \LaTeX{} file by a committed script, which
reproduces its sections, 19 tables, 2 figures, 20 numbered equations and
10 references and fails if its numbering does not match the compiled document.

\section*{Closing}

The revised manuscript, the marked-up manuscript, the supplementary material in
both its \LaTeX{} and Word forms and in both clean and marked versions, and the
point-by-point response accompany this letter. We
remain available to the Editors for the next step of the review process.

\vspace{1.0em}

\begin{minipage}[t]{0.52\linewidth}
\rule{6.2cm}{0.4pt}\\[2pt]
\textbf{Dr. BENIAICHE Ahmed}\\
Corresponding author\\
Laboratory of Fluid Mechanics\\
Ecole Militaire Polytechnique, Algiers, Algeria\\
\href{mailto:beniaiche_ahmed25@yahoo.fr}{beniaiche\_ahmed25@yahoo.fr}
\end{minipage}\hfill
\begin{minipage}[t]{0.42\linewidth}
\raggedleft
\textit{Signed on behalf of all the authors}\\[4pt]
Beniaiche Ahmed\\
Larabi Abderrahim\\
Belkadi Mbarek\\[6pt]
Algiers, 10 November 2026
\end{minipage}

\end{document}
